% Generated draft: requires mathematical and source review.
\documentclass[11pt]{article}
\usepackage[margin=25mm]{geometry}
\usepackage{fontspec}
\setmainfont{Latin Modern Roman}
\usepackage{amsmath,amssymb,mathtools,unicode-math}
\setmathfont{Latin Modern Math}
\usepackage{enumitem,needspace,xurl,hyperref}
\setlength{\parindent}{0pt}
\setlength{\parskip}{5pt}
\setlength{\emergencystretch}{3em}
\newcommand{\N}{\mathbb{N}}
\newcommand{\Z}{\mathbb{Z}}
\newcommand{\Q}{\mathbb{Q}}
\newcommand{\R}{\mathbb{R}}
\newcommand{\C}{\mathbb{C}}
\newcommand{\F}{\mathbb{F}}
\newcommand{\A}{\mathbb{A}}
\newcommand{\PP}{\mathbb{P}}
\newcommand{\E}{\mathbb{E}}
\newcommand{\eps}{\varepsilon}
\newcommand{\dd}{\,\mathrm{d}}
\begin{document}
\section*{UnsolvedMath ID: 1100702}
\textbf{Source problem code:} catalogue entry
\subsection{Definitions and mathematical statement}
Start with a finitely presented group $G$. A graph-of-groups decomposition of a group $H$ means an isomorphism $H\cong\pi_1(\mathcal G)$ with injective edge maps. An edge group is two-ended if it is an infinite finitely generated group with two ends, equivalently a virtually infinite cyclic group; this includes the infinite dihedral possibility and is not restricted to infinite cyclic groups.

The source's alternating hierarchy replaces each current vertex group by a maximal decomposition over finite edge groups, then replaces the resulting vertex groups by maximal decompositions over two-ended edge groups, and repeats these two stages. Maximality is intended to rule out further admissible refinement of the relevant kind. The source does not fix the exact refinement relation, reduction conventions, or relative conditions on incident edge groups. A terminal group admits no further nontrivial decomposition of either prescribed kind; simply choosing its trivial decomposition at one stage does not by itself mean that the entire hierarchy has ended.

The two proposed assertions are:
\begin{enumerate}
\item There exists a finitely presented group $G$ for which this alternating process never terminates.
\item Strong accessibility: for every hyperbolic group $G$, and separately for every CAT(0) group $G$, this alternating process terminates.
\end{enumerate}
Here hyperbolic has the usual finitely generated Gromov-hyperbolic meaning, and a CAT(0) group conventionally admits a proper cocompact isometric action on a proper geodesic CAT(0) space. A fully quantified version must specify whether the first assertion requires some admissible choices or all choices to give nontermination, and whether the second requires termination for every allowed hierarchy or only for an appropriate choice. It must also specify whether termination means finite depth, a finite hierarchy altogether, or absence of an infinite branch. Under finite branching, absence of infinite branches and finiteness of the hierarchy agree, but finite branching is not stated in the record. These choices are not silently resolved here.

The supplied research summary reports termination for hyperbolic groups without $2$-torsion; virtually $2$-torsion-free hyperbolic groups, including residually finite hyperbolic groups; finitely presented subgroups of $\operatorname{SL}(n,\Z)$; Coxeter groups under the stated restriction to minimal splittings; and two-generated torsion-free hyperbolic groups. It attributes these to the cited literature rather than supplying complete hypotheses or proofs. It also mentions a repaired foundation for the earlier hyperbolic argument.

The summary describes Conjecture 1 as open in general, including a torsion-free finitely presented setting with $BS(1,-1)$-type slender subgroups having dihedral quotients; it describes Conjecture 2 as open for arbitrary-torsion hyperbolic groups and for CAT(0) groups in general. A slender group means one whose every subgroup is finitely generated. The phrase ``no noncentral involution'' is identified there as a boundary in the cited work, but its precise scope is not provided. The summary's final sentence is truncated. These are supplied partial-progress assessments, not a newly verified status determination.
\subsection{Short English statement}
Repeatedly split every vertex group, alternating finite and two-ended edge groups. The source conjectures that some finitely presented group allows endless splitting, while hyperbolic and CAT(0) groups always stop. The supplied summary reports many terminating subclasses but leaves the general assertions unresolved. The exact maximality, hierarchy-choice quantifiers, and termination convention are missing and must be fixed before the conjectures become unambiguous.
\subsection{Sources}
\begin{itemize}
\item Local source registry: \url{data/unsolvedmath/problems.json}, permanent UnsolvedMath ID $1100702$, source code catalogue entry, including the supplied PARTIAL-PROGRESS summary. Dataset attribution: Ulam AI and UnsolvedMath Contributors; original-source attribution and rights are retained.
\item Mladen Bestvina, \emph{Questions in Geometric Group Theory}, source-list date 2004, Question 7.2, attributed to Swarup; supplied metadata identifies PDF page 15: \url{https://www.math.utah.edu/~bestvina/eprints/questions-updated.pdf}. Manual transcription and July 29, 2026 access are source metadata; current status and rights remain flagged for review.
\item Vavrichek (2008), the foundation labeled DP01, and Louder--Touikan (2017), Theorem 2.9 and Corollaries 2.7--2.8, as cited by the supplied summary for hyperbolic and linear-group results and the repair of an earlier foundation. Full bibliographic details and exact theorem hypotheses are not supplied.
\item Mihalik--Tschantz, arXiv:1003.0027: \url{https://arxiv.org/abs/1003.0027}, cited for Coxeter groups over minimal splittings; the summary reports an arXiv API check. Kapovich--Weidmann, as cited in Louder--Touikan (2017), is named for the two-generated torsion-free hyperbolic case. No API or linked source was consulted for this draft.
\end{itemize}
\end{document}
